\documentclass[conference]{IEEEtran}
\usepackage{cite}
\usepackage{amsmath,amssymb,amsfonts}
\usepackage{algorithmic}
\usepackage{graphicx}
\usepackage{textcomp}
\usepackage{xcolor}
\def\BibTeX{{\rm B\kern-.05em{\sc i\kern-.025em b}\kern-.08em
    T\kern-.1667em\lower.7ex\hbox{E}\kern-.125emX}}
\begin{document}

\title{Multi-Signal Correlation for Autonomous Remediation in Cloud-Native Microservices: A MAPE-K Driven Approach\\}

\author{\IEEEauthorblockN{1\textsuperscript{st} Given Name Surname}
\IEEEauthorblockA{\textit{dept. name of organization (of Aff.)} \\
\textit{name of organization (of Aff.)}\\
City, Country \\
email address or ORCID}
\and
\IEEEauthorblockN{2\textsuperscript{nd} Given Name Surname}
\IEEEauthorblockA{\textit{dept. name of organization (of Aff.)} \\
\textit{name of organization (of Aff.)}\\
City, Country \\
email address or ORCID}
}

\maketitle

\begin{abstract}
As cloud-native architectures grow in complexity, the demand for autonomous system management becomes critical. Traditional monitoring systems often rely on binary threshold-based alerts, which lack the context necessary to differentiate between transient glitches, resource exhaustion, and security threats. This paper introduces a V3 Self-Healing AI framework designed for high-fidelity microservice environments. By utilizing a Multi-Signal Correlation Engine within a MAPE-K (Monitor-Analyze-Plan-Execute-Knowledge) loop, our system autonomously classifies and remediates system faults. Experimental results demonstrate that the proposed architecture successfully resolves 5xx error spikes, handles service deadlocks, and prevents "Restart Storms" during resource overloads. The framework achieves an average recovery time of under 90 seconds with zero-human intervention.
\end{abstract}

\begin{IEEEkeywords}
Self-Healing AI, Microservices, Prometheus, MAPE-K, Autonomous Remediation, Cloud Computing.
\end{IEEEkeywords}

\section{Introduction}
The evolution of software engineering toward microservice-based architectures has brought unparalleled scalability but introduced significant operational complexity. Ensuring high availability in these distributed systems remains a formidable challenge. Manual intervention is often too slow and prone to human error, specially during large-scale outages.

Existing solutions like Kubernetes liveness probes provide a basic "self-healing" mechanism by restarting unhealthy containers. However, such binary health-checks are blind to "grey failures"---scenarios where a service remains "Running" in a container but is functionally unresponsive due to application-level deadlocks or external dependency timeouts. Furthermore, blind restarts during a resource overload (e.g., a "Flash Sale") can exacerbate system instability, leading to cascading failures.

This paper proposes an intelligent remediation engine that correlates multiple signals---specifically application-level HTTP status codes and infrastructure-level resource utilization (CPU/Memory)---to make precise, context-aware decisions.

\section{Proposed Methodology}
The heart of our system is an autonomous MAPE-K loop integrated with a high-fidelity monitoring stack.

\subsection{Architecture Overview}
The framework consists of four primary modules:
\subsubsection{Monitored Entity} A professional apparel boutique microservice built as a target for autonomous healing.
\subsubsection{Monitoring Layer} A Prometheus-based time-series aggregator that scrapes high-resolution telemetry every 15 seconds.
\subsubsection{Inference & Decision Engine} An AI-driven controller that queries Prometheus and performs cross-metric classification.
\subsubsection{Execution Agent} A Docker-socket integrated handler that performs atomic system operations (Restarts/Scaling).

\subsection{Multi-Signal Correlation Logic}
The system differentiates between two critically different scenarios that both return a 503 Service Unavailable code:
\begin{itemize}
    \item \textbf{Status 503 + Low CPU}: Classified as a \textit{Service Deadlock}. The fix is an immediate restart.
    \item \textbf{Status 503 + High CPU}: Classified as a \textit{Resource Overload}. The AI \textit{suppresses} the restart to preserve active user sessions and prevent a "Restart Storm."
\end{itemize}

\section{Results and Performance Analysis}
The system was evaluated against five major error classes. Table I summarizes the results of the autonomous detection and remediation phases.

\begin{table}[htbp]
\caption{Remediation Precision Matrix}
\begin{center}
\begin{tabular}{|l|c|c|r|}
\hline
\textbf{Scenario}&\textbf{Signal Profile}&\textbf{Action}&\textbf{Heal Time} \\
\hline
Checkout Crash & 500 Spike & Restart & 42s \\
App Deadlock & 503 + Low CPU & Restart & 65s \\
Flash Overload & 503 + High CPU & Suppress & N/A \\
Payment Hang & 504 Spike & Restart & 88s \\
Bot Attack & 429 Spike & Log Only & N/A \\
\hline
\end{tabular}
\end{center}
\end{table}

The experimental data indicates that by using a 60-second observation window, the AI achieves 100\% classification accuracy across the evaluation set, significantly reducing MTTR (Mean Time To Recovery) compared to manual SRE intervention.

\section{Conclusion}
This research demonstrates that multi-signal correlation is vital for secondary-level autonomous healing in cloud-native environments. Our V3 Self-Healing AI framework provides a blueprint for building "resilient-by-design" systems that are aware of their environmental context. Future work will explore the use of Reinforcement Learning to optimize the observation-to-action latency further.

\section*{References}
\begin{itemize}
    \item [1] J. Kephart and D. Chess, ``The vision of autonomic computing,'' in \textit{Computer}, vol. 36, no. 1, 2003.
    \item [2] IBM Corp, ``An architectural blueprint for autonomic computing,'' white paper, 2006.
\end{itemize}

\end{document}
